% §IX --- Conclusion (~0.25 page, ~250 words). Drafted D11.
\section{Conclusion}\label{sec:concl}

We presented \method, a multi-agent hybrid PPO algorithm that jointly
optimizes UAV trajectories, transmit powers, channel selection, and
semantic-symbol cardinalities for multi-UAV DeepSC-VQA semantic
offloading to an MEC server. Four contributions follow. We extended
the conference single-UAV / static-channel formulation to a Dec-POMDP
with multiple UAVs, mobile users, Poisson task arrival, and
time-varying air-to-ground fading. We designed \method as a
centralized-training, decentralized-execution actor with role-aware
specialized discrete heads and a single Gaussian continuous head, and
trained it with a piecewise reward whose four branches couple
communication-resource decisions to LUT-grade semantic-task fidelity.
We derived an additive policy-gradient decomposition and a monotonic
improvement bound that lets the standard PPO clipping mechanism apply
factor-wise. Empirically, \method dominated five baselines on a
measurement-grade DeepSC-VQA benchmark across joint cost, delay,
energy, similarity, and success-rate axes, and held the smallest
degradation slope under multi-UAV scaling, channel and UE-count
perturbations, adversarial jamming, and large-scale UE generalization.

Three directions extend this work. (i) Heterogeneous UAV fleets where
different platforms operate at different altitudes and possess
different sensing modalities; the role-aware actor head pattern is a
natural fit. (ii) Cross-task and cross-dataset generalization, in
which the policy is trained on one DeepSC-VQA LUT and evaluated on
another, possibly online. (iii) Online adaptation via meta-RL or
in-context learning so that a deployed UAV fleet refines its policy
under live channel conditions without resetting the base parameters,
closing the gap between offline-trained semantic encoders and online
adaptive resource allocation.
